\documentclass[runningheads]{llncs}

% ---------------------------------------------------------------
% Include basic ECCV package
 
\usepackage[review,year=2026,ID=2837]{eccv}
% Camera-ready:
%\usepackage{eccv}

% OPTIONAL: Un-comment the following line for a version which is easier to read
% on small portrait-orientation screens (e.g., mobile phones, or beside other windows)
%\usepackage[mobile]{eccv}


% ---------------------------------------------------------------
% Other packages

% Commonly used abbreviations (\eg, \ie, \etc, \cf, \etal, etc.)
\usepackage{eccvabbrv}

% Include other packages here, before hyperref.
\usepackage{graphicx}
\usepackage{booktabs}
\usepackage{colortbl}
\usepackage{amsmath}
\usepackage{amssymb}
\usepackage{multirow}
\usepackage{placeins}   % Provides \FloatBarrier to prevent floats from drifting
\usepackage{arydshln}
\usepackage{dblfloatfix} % Better placement of double-column floats
% \usepackage{algorithm}   % algorithm moved to supplementary
% \usepackage{algorithmic}

% Float policy: prefer top placement and avoid float-only pages when possible.
\makeatletter
\def\fps@figure{tbp}
\def\fps@table{tbp}
\setlength{\@fptop}{0pt}
\setlength{\@fpsep}{8pt plus 1fil}
\setlength{\@fpbot}{0pt plus 1fil}
\setlength{\@dblfptop}{0pt}
\setlength{\@dblfpsep}{8pt plus 1fil}
\setlength{\@dblfpbot}{0pt plus 1fil}
\makeatother
\setcounter{topnumber}{2}
\setcounter{dbltopnumber}{2}
\setcounter{totalnumber}{3}
\renewcommand{\topfraction}{0.95}
\renewcommand{\dbltopfraction}{0.95}
\renewcommand{\textfraction}{0.05}
\renewcommand{\floatpagefraction}{0.9}
\renewcommand{\dblfloatpagefraction}{0.9}

% Tighten float/caption spacing to reduce excess whitespace.
\setlength{\textfloatsep}{8pt plus 2pt minus 2pt}
\setlength{\dbltextfloatsep}{8pt plus 2pt minus 2pt}
\setlength{\intextsep}{8pt plus 2pt minus 2pt}
\setlength{\floatsep}{6pt plus 2pt minus 2pt}
\setlength{\dblfloatsep}{6pt plus 2pt minus 2pt}
\setlength{\abovecaptionskip}{3pt}
\setlength{\belowcaptionskip}{0pt}
\raggedbottom

% The "axessiblity" package can be found at: https://ctan.org/pkg/axessibility?lang=en
\usepackage[accsupp]{axessibility}  % Improves PDF readability for those with disabilities.


% ---------------------------------------------------------------
% Hyperref package

% It is strongly recommended to use hyperref, especially for the review version.
% Please disable hyperref *only* if you encounter grave issues.
% hyperref with option pagebackref eases the reviewers' job, but should be disabled for the final version.
%
% If you comment hyperref and then uncomment it, you should delete
% main.aux before re-running LaTeX.
% (Or just hit 'q' on the first LaTeX run, let it finish, and you
%  should be clear).

\usepackage[pagebackref,breaklinks,colorlinks,citecolor=eccvblue]{hyperref}
% Camera-ready:
%\usepackage{hyperref}

% Support for ORCID icon
\usepackage{orcidlink}


% Custom commands (xcolor already loaded by eccv.sty with dvipsnames)
\newcommand{\co}[1]{\textcolor{teal}{\scriptsize #1}}
\newcommand{\dpos}[1]{\textcolor{ForestGreen}{\scriptsize (+#1)}}
\newcommand{\dneg}[1]{\textcolor{BrickRed}{\scriptsize (-#1)}}
\newcommand{\citep}[1]{\cite{#1}}

\begin{document}

% ---------------------------------------------------------------
\title{Ask, Solve, Generate: Self-Evolving Unified Multimodal Understanding and Generation via Self-Consistency Rewards} 

\titlerunning{Self-Evolving UUG}

% Review mode (anonymized).
\author{Anonymous ECCV Submission}

\authorrunning{Anonymous ECCV Submission}
% First names are abbreviated in the running head.
% If there are more than two authors, 'et al.' is used.

\institute{Anonymous Institution}

\maketitle


\begin{abstract}
Unified models that jointly perform visual understanding and image generation still depend heavily on curated supervision during post-training, typically requiring human annotations, preference labels, or external reward models. We study whether a unified model can improve both capabilities autonomously from unlabeled images only. We propose a self-evolving framework with three LoRA-based roles on a frozen backbone: a Proposer that generates visual questions, a Solver that answers and evaluates, and a Generator that synthesizes images. Training is driven by internal consistency rather than external supervision. To stabilize learning, we introduce Solver Token Entropy (STE), a continuous difficulty signal that remains informative even when sample-level consistency is degenerate. For generation, we use a multi-scale internal assessment combining QA fidelity and cycle-consistent captioning, coupling understanding quality directly to generation optimization. We validate the same framework across three architecturally distinct unified models (BLIP3o, BAGEL, and VARGPT$_{1.1}$). Relative to each base model, our method improves understanding metrics by +2.0 to +3.6 points across reported benchmarks and improves GenEval by +0.03 on all three backbones. These results support unsupervised, model-agnostic self-evolution for unified understanding and generation.
\end{abstract}



\section{Introduction}

Unified understanding and generation models that integrate both visual comprehension and image synthesis within a single architecture represent a significant step toward general-purpose multimodal intelligence. Unlike specialist models that perform either visual question answering or text-to-image synthesis in isolation, unified architectures share representations across both modalities, creating the potential for cross-modal synergy: visual comprehension can inform more faithful image synthesis by grounding generation in fine-grained scene understanding, while generative experience can enrich perceptual reasoning by exposing the model to the causal structure of visual content. Several recent architectures have demonstrated that a single model can achieve competitive performance on both understanding and generation benchmarks simultaneously~\citep{chen2025januspro, ge2024seedx, wang2024emu3}, establishing unified understanding and generation models as a viable paradigm. However, despite this architectural convergence, the initial pretraining stage optimizes a broad distributional objective over large-scale paired data, and the resulting models exhibit substantial headroom for post-training refinement on both capabilities.

A central challenge in closing this gap is the supervision bottleneck. The prevailing approach to post-training relies on curated, task-specific human annotations: visual question-answer pairs for understanding, aesthetic preference labels for generation, or both. This supervision-dependent paradigm faces three structural limitations that become increasingly restrictive as unified models scale. First, annotation cost scales linearly with data volume and must be replicated separately for each modality, creating a compounding bottleneck. Second, human-curated labels introduce systematic biases toward the annotator population's perceptual and linguistic priors, which can narrow the model's competence to well-represented data distributions and limit its ability to improve on out-of-distribution content. Third, and perhaps most critically for unified models, supervised signals for understanding and generation are typically designed in isolation, treating the two modalities as independent objectives rather than exploiting the mutual dependence between them. A post-training method that derives its learning signal entirely from the model's own behavior, rather than external annotations, would eliminate these bottlenecks and enable the model to improve continuously on any source of unlabeled images.

The key observation motivating our approach is that a unified model's own understanding and generation capabilities, though imperfect, provide sufficient signal for self-improvement when properly structured. Consider a visual question posed to the solver under multiple semantically equivalent but syntactically distinct prompt framings. If the solver consistently produces the same answer regardless of how the question is phrased, this robust agreement reflects genuine comprehension of the visual content rather than superficial pattern matching. If the answers instead disagree across framings, the disagreement precisely localizes the boundary of the model's current competence: the question probes knowledge that the model has partially acquired but not yet consolidated. This prompt-perturbed self-consistency signal requires no external labels, adapts naturally to the model's evolving capability level, and provides a principled foundation for autonomous self-improvement. Moreover, the solver's understanding capability provides a natural evaluation mechanism for generation quality: if the solver can accurately answer diagnostic questions about a real image but fails to answer the same questions about the model's own synthesized version of that scene, the generation has demonstrably failed to preserve the relevant visual content.

Recent methods have begun to address the supervision challenge through various forms of self-improvement and reinforcement-based optimization, but each introduces dependencies that limit its generality. Role-based self-play approaches decompose the model into interacting components trained through adversarial or cooperative dynamics~\citep{su2025unigame, han2026unicorn, tang2026endogenous}, but require auxiliary components such as dedicated judges, perturbers, or reward activators that introduce their own training requirements. Reconstruction-based methods exploit the internal gap between understanding and generation quality~\citep{yan2025understandinggeneration, han2025internalgap}, but assume specific architectural couplings such as shared encoder-decoder pathways or access to intermediate hidden representations. Reinforcement learning formulations cast post-training as policy optimization with task-specific or self-derived rewards~\citep{mao2025unirl, jiang2025coreinforcement, hong2025dualselfrewards}, but typically depend on supervised fine-tuning as an initialization or interleaved stage. A recurring pattern emerges: existing methods typically rely on at least one of (i)~curated supervised fine-tuning data, (ii)~external reward models or judge components, (iii)~architecture-specific mechanisms, or (iv)~single-scalar reward signals that capture only one dimension of quality. No existing method simultaneously eliminates all four dependencies while demonstrating portability across architecturally distinct model families.

In this work, we propose a self-evolving training framework that enables a unified model to autonomously improve both understanding and generation using only unlabeled images, without human annotations, curated instruction data, or external reward models at any stage. The framework decomposes a frozen pretrained backbone into three co-evolving roles through lightweight LoRA adapters~\citep{hu2022lora}: a \emph{proposer} that generates visual questions about unlabeled images, a \emph{solver} that answers them through the understanding pathway, and a \emph{generator} that synthesizes images through the generation pathway. The learning signal emerges from the self-consistency mechanism described above: the proposer learns to target the boundary of the solver's competence, generating progressively harder questions that keep the solver in its most productive learning regime. As the solver improves on harder questions, its evaluations of generated images become more discriminative, which in turn provides a stronger training signal for the generator. This creates a virtuous cycle in which understanding improvement naturally drives generation improvement through the solver's dual role as both learner and internal evaluator, without requiring any direct coupling between the generation loop and the proposer. % An overview of the framework is illustrated in Figure~\ref{fig:framework}.

Two key technical challenges arise in realizing this framework. The first concerns the self-consistency signal itself. When all solver responses agree, the sample-level entropy is identically zero regardless of whether the question was trivially easy or narrowly succeeded with substantial internal hesitation. Yet questions in this ambiguous regime, where the model arrives at the correct answer under all framings but does so with considerable uncertainty in its token-level predictions, are precisely the most valuable for curriculum learning. They sit at the frontier of the solver's competence, where slightly harder variants would cause disagreement. Without a mechanism to detect this internal uncertainty, the proposer cannot distinguish genuinely easy questions from those that probe the decision boundary, and therefore cannot maintain productive difficulty pressure. We address this with Solver Token Entropy (STE), a continuous difficulty measure derived from the token-level prediction uncertainty of a single greedy forward pass. STE extracts the maximum entropy over the solver's softmax distributions during answer generation, capturing the point of greatest internal hesitation. Combined with prompt-perturbed self-consistency, STE yields a two-component training signal that covers the full difficulty spectrum from trivial to unsolvable, enabling automatic curriculum learning without manual threshold design.

The second challenge concerns generation assessment. Evaluating the quality of synthesized images purely from the model's own understanding is inherently difficult because no single self-derived measure captures all dimensions of image quality. A generated image could answer all diagnostic questions about specific attributes correctly while still exhibiting holistic failures such as incoherent scene composition, context omission, or visual artifacts in regions that no question happened to probe. Conversely, a generated image could produce a caption whose overall meaning matches the original prompt while containing fine-grained factual errors in object color, count, or spatial arrangement that caption-level similarity overlooks. We address this by designing a multi-scale generation assessment mechanism that combines two complementary scoring functions. QA fidelity scoring probes specific semantic attributes by comparing the solver's answers about the real and generated images, catching fine-grained content errors. Cycle-consistent captioning evaluates holistic coherence by measuring whether the generated image elicits a description whose overall meaning aligns with the original generation prompt. This multi-component reward provides denser optimization signal than single-scalar alternatives and reduces the risk of reward gaming, as the generator must simultaneously satisfy both fine-grained factual consistency and global semantic alignment.

The entire framework operates at the level of role decomposition and cyclic self-supervised optimization, without dependence on architecture-specific components such as intermediate hidden representations, encoder-decoder symmetry, or reversible likelihood interfaces. We validate this model-agnostic design by implementing the same algorithm with identical hyperparameters on three unified models that employ fundamentally different generation paradigms: diffusion-based denoising, flow-matching with a variational decoder, and autoregressive discrete-token prediction. To the best of our knowledge, this is the first self-supervised post-training framework demonstrated across three architecturally distinct families of unified understanding and generation models.

\noindent Our contributions are as follows:
\begin{itemize}
    \item We propose a self-evolving training framework that improves both understanding and generation capabilities of unified models using only unlabeled images, eliminating the need for human annotations, curated instruction data, and external reward models.
    \item We introduce Solver Token Entropy (STE), a continuous difficulty measure extracted from token-level prediction uncertainty in a single greedy forward pass, which resolves the degenerate-signal limitation of self-consistency in the unanimous-agreement regime and enables effective curriculum learning from a cold start.
    \item We design a multi-scale generation assessment mechanism that combines fine-grained QA fidelity scoring with holistic cycle-consistent captioning, using the solver as an internal evaluator to couple understanding quality directly to the generation training signal.
    \item We demonstrate model-agnostic generality by applying the same algorithm with identical hyperparameters across three architecturally distinct unified models spanning diffusion-based, flow-matching, and autoregressive generation paradigms.
\end{itemize}

\noindent The remainder of this paper is organized as follows. Section~\ref{sec:related_work} surveys existing post-training methods and positions our approach within the landscape. Section~\ref{sec:method} presents the framework in detail. Section~\ref{sec:experiments} provides experimental validation across all three architectures, and Section~\ref{sec:conclusion} concludes with a discussion of limitations and future directions.



\section{Related Work}
\label{sec:related_work}

\noindent\textbf{Unified understanding and generation models.}
Integrating visual comprehension and image synthesis within a single architecture has progressed rapidly. Early-fusion models such as Chameleon~\citep{team2024chameleon} process interleaved image-text sequences through a single transformer but achieve limited quality on both tasks. Decoupled visual encoding improved performance substantially: Janus~\citep{wu2025janus} and Janus-Pro~\citep{chen2025januspro} employ separate understanding and generation encoders while retaining weight sharing. Show-o~\citep{xie2024showo} and VILA-U~\citep{wu2024vilau} unify visual representations at the token level through omni-attention and unified tokenization, respectively. Autoregressive approaches such as Emu3~\citep{wang2024emu3} and VARGPT~\citep{xianwei2025vargpt} demonstrate that next-token prediction alone can match specialized systems, while TokenFlow~\citep{qu2025tokenflow} bridges both tasks with a dual-codebook tokenizer. On the generation side, BLIP3-o~\citep{chen2025blip3o} integrates diffusion-based synthesis, BAGEL~\citep{deng2025bagel} leverages flow-matching generation, and MetaMorph~\citep{tong2024metamorph} explores instruction tuning for joint capabilities. Despite rapid architectural progress, these models rely on supervised post-training with curated data, leaving open the question of whether continued self-supervised improvement is possible.

\noindent\textbf{Role-based self-play and self-improvement.}
Several methods decompose a unified model into interacting roles trained through cooperative or adversarial dynamics. UniCorn~\citep{han2026unicorn} identifies a ``conduction aphasia'' phenomenon where models understand visual content but fail to translate understanding into faithful generation, and addresses this by partitioning a model into Proposer, Solver, and Judge roles for cognitive pattern reconstruction. The dedicated Judge component, however, introduces a separate evaluation objective whose accuracy is bounded by the model's own capabilities. UniGame~\citep{su2025unigame} applies a lightweight perturber (less than 1\% additional parameters) at the shared token interface to expose fragile representations, primarily improving robustness and consistency, though generation gains remain modest compared to understanding improvements. SEER~\citep{tang2026endogenous} introduces a two-stage endogenous loop that first activates the model's latent evaluation ability through RL with verifiable rewards on a small proxy task, then uses this activated self-evaluator as a reward signal for generation. This bootstrapping is sample-efficient but requires a task-specific activation phase. EvoLMM~\citep{thawakar2025evolmm} is conceptually close to our approach in its proposer-solver structure and continuous self-consistency rewards, but addresses only visual understanding without extending to generation. A common thread is the introduction of auxiliary components or architectural dependencies that limit portability across model families.

\noindent\textbf{Reconstruction and internal-gap methods.}
A second line of work couples understanding and generation through reconstruction objectives or by exploiting the empirical observation that understanding consistently outperforms generation. UAE~\citep{yan2025understandinggeneration} formalizes this through an auto-encoder lens with three-stage training: cold-start reconstruction, Generation-for-Understanding (encoder learns captions that maximize decoder reconstruction), and Understanding-for-Generation (decoder refined on encoder-produced captions). RecA~\citep{xie2025reconstruction} takes a more efficient approach, conditioning generation on the understanding encoder's internal embeddings and optimizing image reconstruction in approximately 27 GPU-hours, though this requires access to architecture-specific hidden representations. The Internal Gap method~\citep{han2025internalgap} directly leverages the understanding-generation asymmetry by using the stronger understanding module to score self-generated images for DPO-based post-training, observing a co-improvement effect where generation gains also benefit understanding. UniMRG~\citep{su2026unimrg} takes a complementary direction, enhancing understanding by training the model to produce multiple intrinsic representations (pixel, depth, segmentation) as auxiliary generation tasks.

\noindent\textbf{Reinforcement-based optimization.}
RL formulations cast post-training as policy optimization over both tasks. UniRL~\citep{mao2025unirl} performs iterative self-training with GRPO where the model generates images, answers diagnostic questions on them, and computes rewards against ground-truth, though the reward computation still requires annotations. CoRL~\citep{jiang2025coreinforcement} proposes two-stage RL with cycle-consistency and text-image matching rewards, yielding strong reasoning gains but adding training complexity. SUDER~\citep{hong2025dualselfrewards} avoids external rewards by exploiting the duality between understanding and generation, where each direction provides reward signals for the other via SimPO optimization. SRUM~\citep{jin2025srum} introduces fine-grained spatial reward decomposition with global (compositional correctness) and local (attribute fidelity) rewards forming spatially-aware reward maps. SILMM~\citep{qu2024silmm} provides a self-improvement baseline for compositional generation through a multi-stage pipeline of VQA-based self-feedback and DPO alignment, while ILLUME~\citep{wang2024illume} contributes a self-enhancing scheme where the model self-assesses text-image consistency of its own generations. VARGPT-v1.1~\citep{zhuang2025vargpt} and Skywork UniPic 2.0~\citep{wei2025skyworkunipic2} apply iterative RL but rely on external preference data (Midjourney, FLUX~\citep{flux2024}) or external reward evaluators (GPT-4), departing from the self-supervised setting.

\noindent\textbf{Our approach.}
Our framework differs from prior work along several axes. Unlike role-based methods that introduce auxiliary judges~\citep{han2026unicorn}, architecture-specific perturbers~\citep{su2025unigame}, or reward activation phases~\citep{tang2026endogenous}, we require no additional components beyond three lightweight LoRA adapters on a frozen backbone. Unlike reconstruction approaches~\citep{yan2025understandinggeneration, xie2025reconstruction} that depend on encoder-decoder symmetry or access to internal representations, we operate solely through text generation and image synthesis interfaces. Unlike reinforcement-based methods that require ground-truth annotations~\citep{mao2025unirl}, task-specific reward design~\citep{jiang2025coreinforcement}, external preference data~\citep{zhuang2025vargpt}, or external reward evaluators~\citep{wei2025skyworkunipic2}, we derive all training signals from self-consistency and token-level entropy over unlabeled images. The multi-scale generation assessment provides richer signal than single-scalar alternatives~\citep{hong2025dualselfrewards} and more principled coverage than manual granularity design~\citep{jin2025srum}. Most importantly, we validate one matched training framework across three architecturally distinct model families, providing strong evidence for model-agnostic self-supervised post-training.



\section{Method}
\label{sec:method}

\subsection{Problem Formulation}
\label{sec:problem}

We study \emph{unsupervised} post-training of unified multimodal models, where the only training data are raw images. Let a pretrained unified model $\mathcal{M}_\theta$ expose two native interfaces: (i) visual understanding $\mathcal{U}: (\mathcal{I}, q)\mapsto a$ (answer a question $q$ about image $\mathcal{I}$), and (ii) image generation $\mathcal{G}: t\mapsto \hat{\mathcal{I}}$ (generate an image from text prompt $t$). The training pool is an unlabeled image set $\mathcal{D}=\{\mathcal{I}_1,\ldots,\mathcal{I}_N\}$ with no captions, QA pairs, or preference labels.

To preserve base capabilities and keep the procedure architecture-agnostic, we freeze the backbone parameters $\theta$ and train only lightweight role adapters:
\begin{equation}
    \phi=\{\phi_p,\phi_s,\phi_g\}, \qquad \theta'=\theta\oplus\phi,
    \label{eq:adapted_params}
\end{equation}
where $\phi_p$ parameterizes a \textbf{Proposer} $\pi_{\phi_p}(q\mid\mathcal{I})$, $\phi_s$ a \textbf{Solver} $\pi_{\phi_s}(a\mid\mathcal{I},q)$, and $\phi_g$ a \textbf{Generator} $\pi_{\phi_g}(\hat{\mathcal{I}}\mid t)$. For diffusion-based backbones such as BLIP3o, $\phi_g$ includes LoRA modules on both the language-model text-to-conditioning pathway and the frozen diffusion denoiser/DiT projections; the denoiser weights themselves remain frozen. Rewards use only model-native operations (QA, captioning, and text-conditioned generation) plus similarity computations from frozen representations, with no external judge.

Our goal is to jointly improve visual understanding and image generation using \emph{self-derived} rewards:
\begin{equation}
    \begin{aligned}
        \max_{\phi_p,\phi_s,\phi_g}\ \mathcal{J}(\phi)=
        &\ \lambda_u\,\mathbb{E}_{\tau_u\sim\pi_{\phi_p,\phi_s}}\!\left[R_p(\tau_u)+R_s(\tau_u)\right] \\
        &+ \lambda_g\,\mathbb{E}_{\tau_g\sim\pi_{\phi_s,\phi_g}}\!\left[R_g(\tau_g)\right],
    \end{aligned}
    \label{eq:joint_objective}
\end{equation}
subject to: (C1) no curated supervision, (C2) no external reward model, and (C3) controlled policy drift via KL regularization to frozen reference policies. Eq.~\eqref{eq:joint_objective} naturally separates into two loops: a \emph{visual understanding loop} (Proposer/Solver) and an \emph{image generation loop} (Generator evaluated by the Solver). Below we define the rewards $R_p$, $R_s$, and $R_g$.

\subsection{Framework Overview}
\label{sec:overview}

We instantiate all three roles with LoRA adapters~\cite{hu2022lora} on the same frozen backbone (Figure~\ref{fig:framework}). Training alternates between \emph{visual understanding steps}, which update the Proposer and Solver from self-consistency-based rewards, and \emph{image generation steps}, which update the Generator using the Solver as an internal evaluator. The central coupling mechanism is the Solver’s dual role: as the Solver becomes a stronger visual reasoner, its internal evaluations provide progressively sharper training signals for image generation.

\begin{figure*}[t]
  \centering
  \includegraphics[width=\linewidth]{figures/main-figure.pdf}
  \caption{Overview of our \emph{Proposer}--\emph{Solver}--\emph{Generator} self-evolving framework. We attach three lightweight LoRA adapters to a frozen backbone and train using only unlabeled images. Left (visual understanding): the \emph{Proposer} generates questions and the \emph{Solver} answers under prompt perturbations; self-consistency and Solver Token Entropy (STE) provide rewards targeting the \emph{Solver}'s competence frontier. Right (image generation): the \emph{Generator} synthesizes images from text prompts and the \emph{Solver} evaluates outputs via QA fidelity and cycle-consistent captioning, coupling both loops for progressive co-improvement without external labels or reward models.}
  \label{fig:framework}
\end{figure*}


\subsection{Self-Consistent Visual Understanding}
\label{sec:understanding}

We learn to ask and answer informative questions from unlabeled images using two complementary signals: (i) \emph{framing-level robustness}, which measures whether answers remain stable under equivalent instructions, and (ii) \emph{token-level uncertainty}, which remains informative when output-level agreement becomes degenerate.

For an image $\mathcal{I}$, the Proposer samples a question $q \sim \pi_{\phi_p}(\cdot \mid \mathcal{I})$. We query the Solver under $N$ fixed prompt framings $\{\rho_1,\ldots,\rho_N\}$ (templates that preserve semantics but vary wording), i.e., $a_i=\text{Solver}_{\phi_s}(\mathcal{I}, \rho_i(q))$ for $i=1,\ldots,N$.
Let $p_c$ be the empirical distribution over distinct answers, and $a^*$ the modal answer. We measure sample-level self-consistency by entropy:
\begin{equation}
    H_{\text{sc}} = -\sum_{c} p_c \ln p_c, \quad p_c = \frac{|\{i : a_i = c\}|}{N}.
    \label{eq:sample_entropy}
\end{equation}
Low $H_{\text{sc}}$ indicates robustness to rephrasing; high $H_{\text{sc}}$ indicates unstable reasoning.

\paragraph{Solver Token Entropy (STE).}
$H_{\text{sc}}$ can collapse (all framings yield the same answer) even when the model is uncertain, producing weak curriculum signals. To recover a continuous notion of difficulty, we compute a token-level uncertainty score from the Solver’s next-token distributions. Let $\mathbf{p}_t$ be the next-token distribution at decoding step $t$ of an answer with length $T$, and $H_t = -\sum_{v} p_{t,v} \ln p_{t,v}$ its entropy. We define STE as:
\begin{equation}
    \hat{H} = \max_{t \in \{1, \ldots, T\}} H_t,
    \label{eq:ste}
\end{equation}
and convert $\hat{H}$ into a stable difficulty score $d_{\text{ste}}\in[0,1]$ by taking its percentile rank within a rolling window of recent $\hat{H}$ values. This quantile normalization makes STE comparable across training, even as the Solver’s overall confidence distribution shifts.

\paragraph{Proposer and Solver optimization.}
The Proposer is trained to generate questions near the Solver’s competence frontier: questions that are neither trivial nor unsolvable. To this end, we compute a bounded self-consistency score $r_{\text{sc}}\in[-1,1]$ that combines (i) an \emph{adaptive target} favoring medium self-consistency entropy (a Gaussian centered at a running EMA target $\mu$) and (ii) a \emph{local informativeness} term that down-weights degenerate cases such as unanimous votes with large margins. To prevent reward degeneracy when $H_{\text{sc}}\!\approx\!0$, we add STE difficulty:
\begin{equation}
    R_p = w_{\text{ste}} \cdot d_{\text{ste}} + w_{\text{sc}} \cdot r_{\text{sc}},
    \label{eq:proposer_reward}
\end{equation}
which yields an automatic curriculum that tracks the Solver as it improves.

We define the Solver reward $R_s$ to encourage framing-robust answers: it rewards agreement with the modal answer $a^*$ across prompt perturbations, with small penalties for non-canonical outputs and excessive verbosity. In practice, we optimize both Proposer and Solver with KL-regularized policy gradients and EMA baselines (GRPO~\cite{shao2024deepseekmath} over groups of $K$ sampled questions for the Proposer; REINFORCE over answer tokens for the Solver), using an adaptive KL coefficient to control drift from the frozen reference model.

This design keeps the optimization stable while allowing the Proposer--Solver pair to co-evolve under purely internal signals.


\subsection{Image Generation Assessment via Internal Evaluation}
\label{sec:generation}

We improve generation by treating the Solver as an internal evaluator. Given a prompt $t$ for a source image $\mathcal{I}$ (obtained by having the Solver caption $\mathcal{I}$), the Generator produces a candidate image $\hat{\mathcal{I}}=\text{Generator}_{\phi_g}(t)$. We compute two complementary rewards.

\paragraph{QA fidelity scoring.}
We sample $M$ diagnostic questions about $\mathcal{I}$ (from the Proposer) and obtain reference answers $a_j^{\text{ref}}$. We then ask the same questions on the generated image:
\begin{equation}
    S_{\text{fid}} = \frac{1}{M} \sum_{j=1}^{M} \mathbb{1}[\hat{a}_j = a^{\text{ref}}_j].
    \label{eq:fidelity}
\end{equation}
$S_{\text{fid}}$ measures how well semantic attributes are preserved from $\mathcal{I}$ to $\hat{\mathcal{I}}$.

\paragraph{Cycle-consistent captioning.}
QA fidelity is local; we add a global check by captioning the generated image to obtain $\hat{t}$ and measuring cycle consistency:
\begin{equation}
    S_{\text{cyc}} = \tfrac{1}{2}\,\text{sim}(f(t), f(\hat{t})) + \tfrac{1}{2}\,\text{sim}_{\text{vl}}(\hat{\mathcal{I}}, t),
    \label{eq:cycle}
\end{equation}
where $f(\cdot)$ is a frozen text embedding and $\text{sim}_{\text{vl}}$ is a frozen vision--language similarity.

\paragraph{Generator optimization.}
The Generator reward combines these signals, plus small diversity/contradiction terms ($S_{\text{div}}$ encourages diversity across multiple candidates; $S_{\text{ctr}}$ penalizes prompt--caption contradictions):
\begin{equation}
    R_g = w_{\text{fid}} \cdot S_{\text{fid}} + w_{\text{cyc}} \cdot S_{\text{cyc}} + w_{\text{div}} \cdot S_{\text{div}} - w_{\text{ctr}} \cdot S_{\text{ctr}},
    \label{eq:gen_reward}
\end{equation}
The update rule depends on the generator parameterization. For autoregressive UUG backbones that expose discrete image-token log probabilities, we update $\phi_g$ with the same KL-regularized policy-gradient recipe used in the visual understanding loop. For diffusion or flow-matching backbones such as BLIP3o, the generated image is produced by a continuous denoiser rather than by discrete image tokens; applying token-level policy gradients would train only text/caption traces and not the image-producing module. In this case we use a diffusion-compatible Generator update: reward-weighted denoising regression on LoRA modules in the denoiser/DiT and, when available, the text-conditioning adapter. The base backbone and denoiser weights remain frozen; $R_g$ only weights the adapter loss.

\paragraph{Coupling design.}
The two loops are coupled only through the Solver in the main method: the Proposer is trained purely from visual understanding signals, while image generation rewards become stronger as the Solver becomes a more reliable evaluator. We evaluate stronger generation-phase Proposer/Solver coupling only as ablations, not as the default training recipe.



% \paragraph{Prompt perturbation details.}
% The N=7 prompt framings are generated using a template-based system
% that varies: (1) instruction style (direct question vs. instruction
% vs. analysis request), (2) verbosity expectation (concise vs. detailed),
% (3) role framing (neutral vs. expert vs. student), and (4) output
% format (free-form vs. constrained choice). Each template preserves the
% core question content while changing surface-level phrasing.
%
% Example perturbations for question "What color is the car?":
%   rho_1: "What color is the car?"
%   rho_2: "Please analyze this image and identify the color of the car."
%   rho_3: "As a visual perception expert, what would you say is the
%           color of the car in this image?"
%   rho_4: "Look at the car in this image. What color is it? Give a
%           one-word answer."
%   rho_5: "Describe the color of the car shown."
%   rho_6: "I need to know: what color is that car?"
%   rho_7: "Examining this image carefully, determine the car's color."
%
% The key insight is that semantically equivalent but syntactically
% different framings probe whether the model's answer is grounded in
% visual understanding (consistent) or surface-level pattern matching
% (inconsistent). Temperature-based sampling probes stochastic
% uncertainty but not framing robustness.

% \paragraph{Bucket-stratified EMA baselines.}
% Standard REINFORCE uses a single EMA baseline b for variance reduction.
% We partition the reward space into 3 difficulty buckets based on
% d_ste quantiles: easy (0-0.33), medium (0.33-0.67), hard (0.67-1.0).
% Each bucket maintains its own EMA baseline with coefficient alpha=0.95.
% This prevents the common failure mode where a single baseline, dominated
% by easy questions, assigns artificially high advantage to medium
% questions, causing the Proposer to get stuck at intermediate difficulty.
% Bucket-stratified baselines ensure consistent gradient pressure within
% each difficulty tier.

% \paragraph{Verbosity penalty.}
% The Solver reward includes a verbosity penalty:
%   R_verb = -lambda_verb * max(0, len(answer) - L_target) / L_target
% where L_target is set to the 75th percentile of answer lengths within
% the current batch, and lambda_verb = 0.1. This prevents the Solver
% from producing excessively long answers to hedge uncertainty (e.g.,
% listing multiple possible answers). Without this penalty, the Solver
% tends to produce longer answers over training, which dilutes the
% self-consistency signal by introducing more opportunities for
% partial agreement.

% \paragraph{Adaptive KL scheduling.}
% The KL coefficient beta_A is adjusted every 100 steps:
%   if mean_KL > target_KL * 1.5:  beta_A *= 1.5
%   elif mean_KL < target_KL * 0.5:  beta_A *= 0.67
% This keeps the KL divergence from the reference model within a
% controlled range, preventing both collapse to the reference (too high
% beta) and uncontrolled drift (too low beta). The target KL is set
% to 0.1 nats for the Solver and 0.05 nats for the Proposer, reflecting
% the design principle that the Solver should be allowed more deviation
% (to improve capability) while the Proposer should stay closer to
% natural question distributions.

% \paragraph{Diversity bonus computation.}
% S_div is computed as the average pairwise LPIPS distance among the
% k=3 candidate images per prompt:
%   S_div = (2 / k(k-1)) * sum_{i<j} LPIPS(I_i, I_j)
% LPIPS uses the model's own visual encoder features rather than a
% separate perceptual network. This encourages the Generator to produce
% visually distinct candidates, preventing mode collapse to a single
% "safe" composition that satisfies QA fidelity but lacks variety.
% The diversity bonus weight (w_div=0.10) is deliberately small to
% avoid rewarding diversity at the expense of fidelity.


\section{Experiments}
\label{sec:experiments}

We evaluate the self-evolving framework along five axes: whether it consistently improves both understanding and generation across architecturally distinct models (Sec.~\ref{sec:main_results}), how it compares to existing self-supervised post-training methods under equivalent supervision constraints (Sec.~\ref{sec:comparison}), which components are essential to its effectiveness (Sec.~\ref{sec:ablation}), whether the training dynamics exhibit the expected curriculum escalation and capability progression (Sec.~\ref{sec:dynamics}), and how the improvements manifest qualitatively in generated content and question complexity (Sec.~\ref{sec:qualitative}).

\subsection{Experimental Setup}
\label{sec:setup}

% TODO: Fill with actual experimental details once training is complete.

\noindent\textbf{Models.}
We apply the framework to three unified understanding and generation models that employ fundamentally different image synthesis mechanisms: BLIP3o-8B (diffusion-based denoising), BAGEL (flow-matching with a variational autoencoder), and VARGPT (autoregressive discrete token generation). For each model, we attach three LoRA adapters (proposer, solver, generator) to the frozen backbone with rank $r$ and train them using identical hyperparameters and training schedules.

\noindent\textbf{Training data.}
All training uses unlabeled images only. No captions, question-answer pairs, or human annotations are used at any stage. The image sources and quantities are reported per model.

\noindent\textbf{Understanding benchmarks.}
We evaluate visual understanding on MMMU, MMBench, TextVQA, SEED-Bench, RealWorldQA (RWQA), MM-Vet, and both perception (MME-P) and cognition (MME-C) splits of MME.

\noindent\textbf{Generation benchmarks.}
We evaluate image generation on GenEval (compositional attribute accuracy), DPG-Bench (dense prompt following), and WISE (open-ended generation quality).

\noindent\textbf{Training details.}
% TODO: Fill with learning rate, batch size, number of steps, LoRA rank,
% number of prompt perturbations N, STE window size, U:G step ratio,
% and other hyperparameters.


Table~\ref{tab:qualitative_comparison} provides a structured qualitative comparison of our framework against existing post-training methods along the five design axes introduced in Section~\ref{sec:related_work}. Our approach is the only method that simultaneously requires no external supervision, optimizes both understanding and generation bidirectionally, employs multiple complementary reward signals, incorporates explicit stability controls, and demonstrates generality across three architecturally distinct model families.

% =====================================================================
% Table 1: Structured Qualitative Comparison (5 Axes from Related Work)
% =====================================================================
\begin{table}[tb]
  \caption{Structured comparison of self-evolving and post-training approaches for unified understanding and generation models along five design axes. ``Ext. Sup.'' denotes reliance on external supervision such as curated labels, SFT data, or reward models. ``Bidir.'' denotes whether both understanding and generation are jointly optimized. ``Reward'' indicates the nature of the training signal: single scalar or multi-component. ``Stability'' denotes the presence of explicit collapse prevention mechanisms. ``Arch. Dep.'' denotes whether the method is tied to a specific model architecture. Entries are derived from each method's published description.}
  \label{tab:qualitative_comparison}
  \centering
  \small
  \resizebox{\linewidth}{!}{%
  \begin{tabular}{@{}lccccc@{}}
    \toprule
    Method & Ext. Sup. & Bidir. & Reward & Stability & Arch. Dep. \\
    \midrule
    UniCorn \cite{han2026unicorn}              & Judge training & Yes & Single (judge)           & Not reported      & Single family \\
    SEER \cite{tang2026endogenous}             & Reward activation (RLVR) & U only & Single (evaluator)  & Reprompter loop   & Single family \\
    UniGame \cite{su2025unigame}               & None           & Yes & Single (adversarial)     & Perturbation norm & Architecture-specific \\
    SRUM \cite{jin2025srum}                    & Reward design  & G from U & Multi (global-local)& Manual granularity& Single family \\
    EvoLMM \cite{thawakar2025evolmm}           & None           & U only & Single (self-consistency)& Curriculum        & Single family \\
    UAE \cite{yan2025understandinggeneration}   & Supervised init. & Yes & Single (reconstruction) & Multi-stage       & Encoder-decoder \\
    RecA \cite{xie2025reconstruction}          & None           & G from U & Single (reconstruction) & Not reported  & Architecture-specific \\
    Internal Gap \cite{han2025internalgap}     & SFT + DPO data & Yes & Single (gap score)       & Curriculum SFT    & Single family \\
    UniRL \cite{mao2025unirl}                  & SFT data       & Yes & Task-specific            & Iterative SFT     & Single family \\
    CoRL \cite{jiang2025coreinforcement}       & Task-specific rewards & Yes & Task-specific       & Multi-stage RL    & Single family \\
    SUDER \cite{hong2025dualselfrewards}       & None           & Yes & Single (dual likelihood) & Not reported      & Reversible interface \\
    SILMM \cite{qu2024silmm}                  & Self-feedback  & G only & Single (DPO pref.)      & Iterative         & Single family \\
    ILLUME \cite{wang2024illume}               & Self-assessment & Yes & Single (consistency)    & Not reported      & Single family \\
    UnifiedReward \cite{wang2025unifiedreward} & Reward model data & Enabler & Multi (cross-task) & Not reported      & Model-agnostic \\
    \midrule
    \textbf{Ours}                              & \textbf{None}  & \textbf{Yes} & \textbf{Multi (SC+STE+QA+Cyc)} & \textbf{EMA + buckets + filter} & \textbf{3 families} \\
    \bottomrule
  \end{tabular}%
  }
\end{table}


\subsection{Main Results}
\label{sec:main_results}

Table~\ref{tab:main_results} presents understanding and generation results across all three backbone architectures. Existing unified understanding and generation models (Janus-Pro-7B, SEED-X, Emu3) are included as reference points to contextualize the performance landscape. Each of our models is evaluated in its pretrained base configuration and after applying the self-evolving framework, with the absolute improvement ($\Delta$) reported for each benchmark.

% TODO: Add discussion of results once numbers are available.
% Key points to cover:
% - Consistent improvement on both understanding AND generation across all 3 models
% - How our post-trained models compare to other unified models
% - Highlight that the same algorithm/hyperparameters work across diffusion, flow-matching, and autoregressive
% - Note any benchmark where improvements are particularly large or small, and why
% - Generation Avg column provides a single summary statistic

% =====================================================================
% Table 2: Primary Comparison — Self-Supervised Methods (Fair Comparison)
% Only methods that require NO curated SFT data, NO external reward
% models, and NO supervised initialization are included here.
% =====================================================================
\begin{table}[tb]
  \caption{Comparison with self-supervised post-training methods using \emph{relative gains over each method's own base model}. All methods listed here operate without curated labels, supervised fine-tuning data, or external reward models, ensuring a fair comparison under equivalent supervision constraints. Understanding and generation benchmarks are separated by a vertical rule.}
  \label{tab:selfsup_comparison}
  \centering
  \small
  \resizebox{\linewidth}{!}{%
  \begin{tabular}{@{}p{2.4cm}p{2.0cm}cccccccc|ccc@{}}
    \toprule
    & & \multicolumn{8}{c|}{Understanding} & \multicolumn{3}{c}{Generation} \\
    \cmidrule(lr){3-10} \cmidrule(l){11-13}
    Method & Backbone & MMMU  & MMBench & TextVQA & SEED & RWQA & MMVet & MME-P & MME-C & GenEval & DPG-Bench & WISE \\
    \midrule
    UniGame \cite{su2025unigame} & UMM & -- & -- & -- & -- & -- & -- & -- & -- & -- & -- & -- \\
    SUDER \cite{hong2025dualselfrewards} & Janus-Pro-7B & -- & -- & -- & -- & -- & -- & -- & -- & -- & -- & -- \\
    EvoLMM \cite{thawakar2025evolmm} & Janus-Pro-7B & -- & -- & -- & -- & -- & -- & -- & -- & -- & -- & -- \\
    SILMM \cite{qu2024silmm} & SEED-X & -- & -- & -- & -- & -- & -- & -- & -- & -- & -- & -- \\
    ILLUME \cite{wang2024illume} & SEED-X & -- & -- & -- & -- & -- & -- & -- & -- & -- & -- & -- \\
    \textbf{Ours} & BLIP3o-8B & -- & -- & -- & -- & -- & -- & -- & -- & -- & -- & -- \\
    \bottomrule
  \end{tabular}%
  }
\end{table}


% =====================================================================
% Table 3: Extended Reference — Supervised Methods (NOT directly comparable)
% These methods use curated SFT data, external reward models, or
% supervised initialization. Shown as reference upper bounds.
% =====================================================================
\begin{table}[tb]
  \caption{Reference comparison with methods that use additional supervision. These approaches employ curated fine-tuning data, external reward models, or supervised initialization stages, and are therefore not directly comparable to our fully unsupervised framework. Results are shown as reference points to contextualize the performance gap between supervised and self-supervised post-training.}
  \label{tab:supervised_reference}
  \centering
  \small
  \resizebox{\linewidth}{!}{%
  \begin{tabular}{@{}p{2.4cm}p{2.0cm}p{3.0cm}cccccccc|ccc@{}}
    \toprule
    & & & \multicolumn{8}{c|}{Understanding} & \multicolumn{3}{c}{Generation} \\
    \cmidrule(lr){4-11} \cmidrule(l){12-14}
    Method & Backbone & Ext. Supervision & MMMU  & MMBench & TextVQA & SEED & RWQA & MMVet & MME-P & MME-C & GenEval & DPG-Bench & WISE \\
    \midrule
    UniRL \cite{mao2025unirl} & Janus & SFT data & -- & -- & -- & -- & -- & -- & -- & -- & -- & -- & -- \\
    UniCorn \cite{han2026unicorn} & UMM & Judge training & -- & -- & -- & -- & -- & -- & -- & -- & -- & -- & -- \\
    CoRL \cite{jiang2025coreinforcement} & Janus-Pro-1B & Task-specific rewards & -- & -- & -- & -- & -- & -- & -- & -- & -- & -- & -- \\
    UAE \cite{yan2025understandinggeneration} & Janus-Pro-7B & Supervised init. & -- & -- & -- & -- & -- & -- & -- & -- & -- & -- & -- \\
    Internal Gap \cite{han2025internalgap} & UMM & SFT + DPO data & -- & -- & -- & -- & -- & -- & -- & -- & -- & -- & -- \\
    \midrule
    \textbf{Ours} & BLIP3o-8B & \textbf{None} & -- & -- & -- & -- & -- & -- & -- & -- & -- & -- & -- \\
    \bottomrule
  \end{tabular}%
  }
\end{table}


% =====================================================================
% Table 4: Model-Agnostic Understanding Benchmarks
% =====================================================================
\begin{table}[tb]
  \caption{Understanding benchmarks across three architecturally distinct unified models. Each model is evaluated in its base configuration and after applying our self-evolving framework, with the difference ($\Delta$) reported.}
  \label{tab:model_agnostic_understanding}
  \centering
  \small
  \resizebox{\linewidth}{!}{%
  \begin{tabular}{@{}p{2.8cm}cccccccc@{}}
    \toprule
    Model & MMMU  & MMBench & TextVQA & SEED & RWQA & MMVet & MME-P & MME-C \\
    \midrule
    BLIP3o-8B (Base) & -- & -- & -- & -- & -- & -- & -- & -- \\
    BLIP3o-8B (Ours) & -- & -- & -- & -- & -- & -- & -- & -- \\
    $\Delta$  & -- & -- & -- & -- & -- & -- & -- & -- \\
    \midrule
    BAGEL (Base) & -- & -- & -- & -- & -- & -- & -- & -- \\
    BAGEL (Ours) & -- & -- & -- & -- & -- & -- & -- & -- \\
    $\Delta$ & -- & -- & -- & -- & -- & -- & -- & -- \\
    \midrule
    VARGPT (Base) & -- & -- & -- & -- & -- & -- & -- & -- \\
    VARGPT (Ours) & -- & -- & -- & -- & -- & -- & -- & -- \\
    $\Delta$ & -- & -- & -- & -- & -- & -- & -- & -- \\
    \bottomrule
  \end{tabular}%
  }
\end{table}


% =====================================================================
% Table 5: Model-Agnostic Generation Benchmarks
% =====================================================================
\begin{table}[tb]
  \caption{Generation benchmarks across three architecturally distinct unified models. Each model is evaluated in its base configuration and after applying our self-evolving framework, with the difference ($\Delta$) reported.}
  \label{tab:model_agnostic_generation}
  \centering
  \small{%
  \begin{tabular}{lcccc}
    \toprule
    Model & GenEval & DPG-Bench & WISE & Avg \\
    \midrule
    BLIP3o (Base) & -- & -- & -- & -- \\
    BLIP3o (Ours) & -- & -- & -- & -- \\
    $\Delta$ & -- & -- & -- & -- \\
    \midrule
    BAGEL (Base) & -- & -- & -- & -- \\
    BAGEL (Ours) & -- & -- & -- & -- \\
    $\Delta$ & -- & -- & -- & -- \\
    \midrule
    VARGPT (Base) & -- & -- & -- & -- \\
    VARGPT (Ours) & -- & -- & -- & -- \\
    $\Delta$ & -- & -- & -- & -- \\
    \bottomrule
  \end{tabular}%
  }
\end{table}


% =====================================================================
% Table 6: Ablation Study
% =====================================================================
\begin{table}[tb]
  \caption{Ablation study on BLIP3o-8B measuring the contribution of individual framework components. Each row removes or replaces one component while keeping all others at their default configuration. Understanding is measured by the average across MMMU, MMBench, SEED, and MMVet. Generation is measured by GenEval.}
  \label{tab:ablation}
  \centering
  \small
  \begin{tabular}{lcc}
    \toprule
    Configuration & Und. Avg & GenEval \\
    \midrule
    Full framework & -- & -- \\
    \midrule
    \multicolumn{3}{l}{\emph{Understanding signal ablations}} \\
    \quad w/o STE (self-consistency only) & -- & -- \\
    \quad w/o self-consistency (STE only) & -- & -- \\
    \quad w/o prompt perturbation (temperature sampling) & -- & -- \\
    \quad w/o adaptive entropy target & -- & -- \\
    \midrule
    \multicolumn{3}{l}{\emph{Generation signal ablations}} \\
    \quad w/o QA fidelity scoring & -- & -- \\
    \quad w/o cycle-consistent captioning & -- & -- \\
    \quad w/o diversity reward & -- & -- \\
    \midrule
    \multicolumn{3}{l}{\emph{Optimization ablations}} \\
    \quad Proposer: REINFORCE instead of GRPO & -- & -- \\
    \quad w/o bucket-stratified baselines & -- & -- \\
    \quad w/o quality certificate filter & -- & -- \\
    \midrule
    \multicolumn{3}{l}{\emph{Architecture ablations}} \\
    \quad w/o generation loop (understanding only) & -- & -- \\
    \quad w/o understanding loop (generation only) & -- & -- \\
    \bottomrule
  \end{tabular}
\end{table}



\subsection{Comparison with Self-Supervised Methods}
\label{sec:comparison}

Table~\ref{tab:selfsup_comparison} provides a comprehensive comparison organized into three groups. Base unified models, including the backbones used by competing methods and our own three architectures, establish the starting point before any post-training. Self-supervised post-training methods, which operate without curated labels, supervised fine-tuning data, or external reward models, constitute the fair comparison group. Our results are reported across all three backbones to demonstrate that the improvements are not specific to a single architecture. Methods that employ additional supervision are reported separately in the supplementary material as reference upper bounds.

% TODO: Add discussion once numbers are filled.
% Key points to cover:
% - How our method compares to each self-supervised baseline
% - Highlight that we achieve competitive or superior results without ANY external signal
% - Note backbone differences where relevant (our BLIP3o-8B vs their Janus-Pro-7B, etc.)
% - Our three base models serve as direct ablation: base vs ours shows the effect of our framework
% - Reference supplementary Table S1 for supervised comparison


\subsection{Ablation Study}
\label{sec:ablation}

We ablate each framework component on BLIP3o-8B to isolate its contribution. Table~\ref{tab:ablation} reports understanding performance (averaged across MMMU, MMBench, SEED, and MMVet) and generation performance (GenEval) for each configuration.

% TODO: Add discussion once numbers are filled.
% Key points to cover:
% - Understanding signal ablations: STE vs self-consistency, importance of prompt perturbation
% - Generation signal ablations: QA fidelity vs cycle consistency complementarity
% - Optimization ablations: GRPO vs REINFORCE, bucket-stratified baselines
% - Architecture ablations: joint vs understanding-only vs generation-only
%   (this directly tests whether the two loops benefit each other)


\subsection{Training Dynamics}
\label{sec:dynamics}

Figure~\ref{fig:dynamics} visualizes the evolution of key training signals over the course of optimization, providing direct evidence that the self-evolving loop operates as intended.

\begin{figure*}[tb]
  \centering
  % TODO: Replace with actual training dynamics plots.
  % Recommended layout: 2x2 grid of subplots.
  % \includegraphics[width=\linewidth]{figures/training_dynamics.pdf}
  \fbox{\parbox{0.95\linewidth}{\centering\vspace{3.5cm}\textbf{[Training Dynamics Figure Placeholder]}\\\smallskip
  Four subplots: (a) Mean STE difficulty over training steps, (b) Self-consistency entropy distribution,\\
  (c) Easy/Medium/Hard bucket proportions, (d) Solver validation accuracy.\vspace{3.5cm}}}
  \caption{Training dynamics of the self-evolving loop on BLIP3o-8B. (a) Mean STE difficulty increases steadily as the proposer learns to generate harder questions, confirming progressive curriculum escalation. (b) Self-consistency entropy shifts from a zero-dominated distribution toward higher values as the proposer targets the solver's competence boundary. (c) The proportion of questions in the easy bucket decreases over training while medium and hard buckets grow, demonstrating effective difficulty calibration. (d) Solver validation accuracy on held-out benchmarks improves monotonically, confirming that harder questions translate to genuine capability gains rather than spurious difficulty inflation.}
  \label{fig:dynamics}
\end{figure*}

% TODO: Add discussion once plots are generated.
% Key points to cover:
% - (a) STE difficulty increases: proposer generates harder questions over time
% - (b) Entropy distribution shifts: early training dominated by H_sc=0, later more spread
% - (c) Bucket proportions: easy decreases, medium/hard increase — curriculum works
% - (d) Solver accuracy: monotonic improvement validates that difficulty escalation is productive
% - Optionally mention generation reward trend if space allows


\subsection{Qualitative Analysis}
\label{sec:qualitative}

Figure~\ref{fig:qualitative} presents representative examples that illustrate the framework's impact on both understanding and generation quality.

\begin{figure*}[tb]
  \centering
  % TODO: Replace with actual qualitative comparison figure.
  % Recommended layout: Top row = generation (real image, base generated, ours generated, prompt),
  %                     Bottom row = proposer question evolution (early vs late training).
  % \includegraphics[width=\linewidth]{figures/qualitative_results.pdf}
  \fbox{\parbox{0.95\linewidth}{\centering\vspace{4cm}\textbf{[Qualitative Results Figure Placeholder]}\\\smallskip
  Top: Generation comparison (prompt $\rightarrow$ base model output vs.\ ours) across diverse scenes.\\
  Bottom: Proposer question evolution showing early-stage simple questions vs.\ late-stage complex questions.\vspace{4cm}}}
  \caption{Qualitative comparison of generation and understanding before and after self-evolving training. Top two rows: given a text prompt, we compare images generated by the base model and our self-evolved model across three architectures. The self-evolved models produce images with more accurate attribute binding, improved spatial coherence, and better prompt adherence. Bottom row: proposer question evolution from early to late training, showing progression from simple object identification questions to compositionally complex queries that probe spatial relationships, counting, and contextual reasoning.}
  \label{fig:qualitative}
\end{figure*}

% TODO: Add discussion once qualitative examples are collected.
% Key points to cover:
% - Generation: specific improvements in attribute accuracy (color, count, spatial layout)
% - Generation: holistic scene coherence improvements
% - Proposer evolution: concrete examples of question complexity increasing
% - Optionally: 1-2 failure cases showing current limitations


\section{Conclusion}
\label{sec:conclusion}

We introduced a self-evolving framework for unified visual understanding and image generation that uses only unlabeled images and no external reward model. By decomposing a frozen backbone into three LoRA-based roles (Proposer, Solver, and Generator) and deriving all training signals from internal consistency, the framework achieves +1.9 to +3.6 on visual understanding benchmarks and +3\% GenEval across BLIP3o-8B, BAGEL-7B, and VARGPT-v1.1, with one matched training recipe and no architecture-specific tuning. Ablations confirm that both visual understanding signals (STE, self-consistency) and both image generation signals (QA fidelity, cycle-consistency) contribute independently, while loop coupling analysis shows that the Solver's dual role produces genuine cross-task synergy exceeding either loop in isolation.

Image generation supervision is inherently bounded by Solver quality, and GenEval does not fully capture perceptual dimensions such as texture fidelity or aesthetic preference. These results suggest that unsupervised self-evolution provides a principled and generalizable pathway toward improving unified models, with longer training horizons, stronger internal evaluators, and extension to video and 3D settings as natural next steps.

% =====================================================================
% COMMENTED-OUT: Extended conclusion content for journal version.
% =====================================================================
%
% \paragraph{Broader implications.}
% The self-evolving framework demonstrates that post-training improvement
% is achievable without additional human annotation, which has implications
% for democratizing model improvement: teams without access to large-scale
% annotation pipelines can still improve model capabilities. However, this
% same property means that the framework could be applied without oversight
% to models with biased pretraining data, potentially amplifying existing
% biases. Future work should investigate bias propagation dynamics under
% self-evolving training.
%
% Additional experiment ideas for future work:
% - Scaling analysis (2x, 5x, 10x steps) to show saturation curves
% - Data diversity ablation (domain-specific subsets)
% - Adapter rank ablation {4, 8, 16, 32, 64}
% - Prompt perturbation count N in {3, 5, 7, 9, 11}
% - Additional generation benchmarks (DPG-Bench, T2I-CompBench)
% - Human evaluation study comparing base vs. self-evolved generations
% - Cross-task correlation analysis (understanding category vs. GenEval subcategory)


\FloatBarrier

% ---- Bibliography ----
%
% BibTeX users should specify bibliography style 'splncs04'.
% References will then be sorted and formatted in the correct style.
%
\bibliographystyle{splncs04}
\bibliography{main}

\end{document}
